\section{Limitations}
\label{sec:limitations}

\paragraph{Instance Scale.}
Our experiments use synthetic instances up to 25 scenes. While this covers the range studied in prior MSSP literature~\citep{liu2019, long2020}, real-world feature films may involve 100+ scenes. The decomposition approach should scale to such sizes (producing more blocks while keeping each block small), but empirical validation on larger instances remains future work.

\paragraph{Synthetic Data.}
All instances are synthetically generated following established parameter ranges. While our generator incorporates precedence, time windows, and geographic cost structures, it may not capture all idiosyncrasies of real production schedules (e.g., weather dependencies, union regulations, equipment sharing). Validation on real film scripts, such as the 70-scene instance used by \citet{tekin2023}, would strengthen external validity.

\paragraph{Block Granularity.}
When Louvain decomposition produces too few blocks ($K \leq 4$), the block-permutation search space is too small for evolutionary search to provide value. On S15 ($K=3$, $3!=6$ permutations) and S20 ($K=4$, $4!=24$ permutations), BDM underperforms scene-level metaheuristics. Hierarchical decomposition or sub-block splitting could address this.

\paragraph{Pareto Front Collapse.}
The greedy decoder assigns identical makespan regardless of scene ordering, causing the Pareto front to degenerate to a single point on all tested instances. The bi-objective formulation does not produce genuine cost-makespan trade-offs under the current decoder. A decoder that allows variable daily utilization would be needed to realize multi-objective benefits.

\paragraph{Number of Seeds.}
We use 3 independent seeds per algorithm. While sufficient for initial comparisons, 30 seeds would provide stronger statistical power, as recommended by~\citet{derrac2011}. Computational budget constraints limited our seed count.
